%-----------------------------------------------------------------------------------------------------------------------------------------------%
%	The MIT License (MIT)
%
%	Copyright (c) 2015 Jan Küster
%
%	Permission is hereby granted, free of charge, to any person obtaining a copy
%	of this software and associated documentation files (the "Software"), to deal
%	in the Software without restriction, including without limitation the rights
%	to use, copy, modify, merge, publish, distribute, sublicense, and/or sell
%	copies of the Software, and to permit persons to whom the Software is
%	furnished to do so, subject to the following conditions:
%	
%	THE SOFTWARE IS PROVIDED "AS IS", WITHOUT WARRANTY OF ANY KIND, EXPRESS OR
%	IMPLIED, INCLUDING BUT NOT LIMITED TO THE WARRANTIES OF MERCHANTABILITY,
%	FITNESS FOR A PARTICULAR PURPOSE AND NONINFRINGEMENT. IN NO EVENT SHALL THE
%	AUTHORS OR COPYRIGHT HOLDERS BE LIABLE FOR ANY CLAIM, DAMAGES OR OTHER
%	LIABILITY, WHETHER IN AN ACTION OF CONTRACT, TORT OR OTHERWISE, ARISING FROM,
%	OUT OF OR IN CONNECTION WITH THE SOFTWARE OR THE USE OR OTHER DEALINGS IN
%	THE SOFTWARE.
%	
%
%-----------------------------------------------------------------------------------------------------------------------------------------------%


%============================================================================%
%
%	DOCUMENT DEFINITION
%
%============================================================================%

%we use article class because we want to fully customize the page and dont use a cv template
\documentclass[10pt,A4]{article}	


%----------------------------------------------------------------------------------------
%	ENCODING
%----------------------------------------------------------------------------------------

%we use utf8 since we want to build from any machine
\usepackage[utf8]{inputenc}		

%----------------------------------------------------------------------------------------
%	LOGIC
%----------------------------------------------------------------------------------------

% provides \isempty test
\usepackage{xifthen}

%----------------------------------------------------------------------------------------
%	FONT
%----------------------------------------------------------------------------------------

% some tex-live fonts - choose your own

%\usepackage[defaultsans]{droidsans}
%\usepackage[default]{comfortaa}
%\usepackage{cmbright}
%\usepackage[default]{raleway}
%\usepackage{fetamont}
%\usepackage[default]{gillius}
%\usepackage[light,math]{iwona}
%\usepackage[thin]{roboto} 

% set font default
\renewcommand*\familydefault{\sfdefault} 	
\usepackage[T1]{fontenc}

% more font size definitions
\usepackage{moresize}		


%----------------------------------------------------------------------------------------
%	PAGE LAYOUT  DEFINITIONS
%----------------------------------------------------------------------------------------

%debug page outer frames
%\usepackage{showframe}			


%define page styles using geometry
\usepackage[a4paper]{geometry}		

% for example, change the margins to 2 inches all round
\geometry{top=20mm, bottom=20mm, left=1.5cm, right=1.5cm} 	

%use customized header
\usepackage{fancyhdr}
\pagestyle{fancy}
\fancyhf{}

%use our custom fancy header definitions
\fancyfoot[R]{\thepage}
\renewcommand{\headrulewidth}{0pt}
\renewcommand{\footrulewidth}{0pt}

%less space between header and content
\setlength{\headheight}{-5pt}		


%customize entries left, center and right
%\lhead{}
%\chead{}
%\rhead{}


%indentation is zero
\setlength{\parindent}{0mm}

%----------------------------------------------------------------------------------------
%	TABLE /ARRAY DEFINITIONS
%---------------------------------------------------------------------------------------- 

%for layouting tables
\usepackage{multicol}			
\usepackage{multirow}

%extended aligning of tabular cells
\usepackage{array}

\newcolumntype{x}[1]{%
>{\raggedleft\hspace{0pt}}p{#1}}%


%----------------------------------------------------------------------------------------
%	GRAPHICS DEFINITIONS
%---------------------------------------------------------------------------------------- 

%for header image
\usepackage{graphicx}

%for floating figures
\usepackage{wrapfig}
\usepackage{float}
%\floatstyle{boxed} 
%\restylefloat{figure}

%for drawing graphics		
\usepackage{tikz}				
\usetikzlibrary{shapes, backgrounds,mindmap, trees}


%----------------------------------------------------------------------------------------
%	Color DEFINITIONS
%---------------------------------------------------------------------------------------- 

%accent color
\definecolor{sectcol}{RGB}{77, 111, 196}

\definecolor{sectcoldark}{RGB}{125, 62, 3}

%dark background color
\definecolor{bgcol}{RGB}{110,110,110}

%light background / accent color
\definecolor{softcol}{RGB}{225,225,225}

% Links
\usepackage{color}
\usepackage{hyperref}
\hypersetup{
    colorlinks=true,
    urlcolor=sectcoldark
    }



%============================================================================%
%
%
%	DEFINITIONS
%
%
%============================================================================%

%----------------------------------------------------------------------------------------
% 	HEADER
%----------------------------------------------------------------------------------------

%remove pagenum
% \renewcommand{\thepage}{}	

%remove section num		
\renewcommand{\thesection}{}			

%----------------------------------------------------------------------------------------
% 	ARROW GRAPHICS in Tikz
%----------------------------------------------------------------------------------------

% a six pointed arrow poiting to the left
\newcommand{\tzlarrow}{(0,0) -- (0.2,0) -- (0.3,0.2) -- (0.2,0.4) -- (0,0.4) -- (0.1,0.2) -- cycle;}	

% include the left arrow into a tikz picture
% param1: fill color
%
\newcommand{\larrow}[1]
{\begin{tikzpicture}[scale=0.58]
	 \filldraw[fill=#1!100,draw=#1!100!black]  \tzlarrow
 \end{tikzpicture}
}

% a six pointed arrow poiting to the right
\newcommand{\tzrarrow}{ (0,0.2) -- (0.1,0) -- (0.3,0) -- (0.2,0.2) -- (0.3,0.4) -- (0.1,0.4) -- cycle;}

% include the right arrow into a tikz picture
% param1: fill color
%
\newcommand{\rarrow}
{\begin{tikzpicture}[scale=0.7]
	\filldraw[fill=sectcol!100,draw=sectcol!100!black] \tzrarrow
 \end{tikzpicture}
}



%----------------------------------------------------------------------------------------
%	custom sections
%----------------------------------------------------------------------------------------

% create a coloured box with arrow and title as cv section headline
% param 1: section title
%
\newcommand{\cvsection}[1]
{
	\begin{center}
		\large\textcolor{sectcol}{\textbf{#1}}
	\end{center}
}

%create a coloured arrow with title as cv meta section section
% param 1: meta section title
%
\newcommand{\metasection}[2]
{
%\begin{tabular*}{1\textwidth}{r r}
\footnotesize{#2} \hspace*{\fill} \footnotesize{#1}\\[1pt]
%\end{tabular*}
}

%----------------------------------------------------------------------------------------
%	 CV EVENT
%----------------------------------------------------------------------------------------

% creates a stretched box as cv entry headline followed by some paragraphs about 
% the work you did
% param 1:	event time i.e. 2014 or 2011-2014 etc.
% param 2:	event name (what did you do?)
% param 3:	institution (where did you work / study)
% param 4:	list of paragraphs
%
\newcommand{\cvevent}[4]
{

\begin{tabular*}{1\textwidth}{p{13.6cm}  x{3.9cm}}
	\textbf{#2} - \textcolor{bgcol}{#3} &   \vspace{2.5pt}\textcolor{sectcol}{#1}
\end{tabular*}

\vspace{-8pt}
\textcolor{softcol}{\hrule}
\vspace{6pt}

	\foreach \desc in {#4}{
		$\cdot$ \desc\\[3pt]
	}
	
\vspace{3pt}
}

% creates a stretched box as 
\newcommand{\cveventmeta}[2]
{
	\mbox{\mystrut \hspace{87pt}\textit{#1}}\\
	#2
}

%----------------------------------------------------------------------------------------
% CUSTOM STRUT FOR EMPTY BOXES
%----------------------------------------- -----------------------------------------------
\newcommand{\mystrut}{\rule[-.3\baselineskip]{0pt}{\baselineskip}}

%============================================================================%
%
%
%
%	DOCUMENT CONTENT
%
%
%
%============================================================================%
\begin{document}



%----------------------------------------------------------------------------------------
%	HEADER IMAGE
%----------------------------------------------------------------------------------------

%\begin{figure}[H]
%\begin{flushright}
%	\includegraphics[trim= 320 130 460 210,clip,width=0.2\linewidth]{myfoto.jpg}	%trimming relative to image size!
%\end{flushright}
%\end{figure}


%---------------------------------------------------------------------------------------
%	TITLE HEADLINE
%----------------------------------------------------------------------------------------
\vspace{-8pt}
\begin{center}
	\HUGE \textsc{Rachel Bolufé Torres} \textcolor{sectcol}{\rule[-1mm]{1mm}{0.9cm}} \textsc{Resume}\\[2pt]
	\small Profesora e Investigadora
\end{center}



\vspace{6pt}


%---------------------------------------------------------------------------------------
%	META SECTION
%----------------------------------------------------------------------------------------
\metasection{Madrid, España}{\textbf{MSc.:} Antropología Física y Evolución Humanas}
\metasection{<<phonerachel>>}{\textbf{Lic.:} Biología}
\metasection{<<emailrachel>>}{\textbf{Idomas:} Spanish (Nativo), Portugués (Profesional), Inglés (Básico)}
\metasection{}{\textbf{Campos:} Investigación, Profesorado} 
\metasection{}{\textbf{Social:} \href{https://www.linkedin.com/in/rachel-boluf\%C3\%A9-torres-5b6611a7/}{Linkedin}}
\vspace{-2pt}
\textcolor{softcol}{\hrule}
\vspace{6pt}

\normalsize

%---------------------------------------------------------------------------------------
%	SUMMARY (optional)
%----------------------------------------------------------------------------------------
\vspace{-6pt}
\cvsection{Resumen}
Soy una profesora e investigadora apasionada por la biología y la antropología física. Me gusta la investigación, el aprendizaje y la enseñanza.


%============================================================================%
%
%	CV SECTIONS AND EVENTS (MAIN CONTENT)
%
%============================================================================%

%---------------------------------------------------------------------------------------
%	EXPERIENCE
%----------------------------------------------------------------------------------------
\cvsection{Experiencia}

\cvevent{Nov 2023 - Present}{Profesora}{Academa Stellar (\href{https://academiastellar.com/}{academiastellar.com})}{
	{Profesora de Refuerzo de ciencias para estudiantes de universidades como Francisco de Vitoria, Europea de Madrid, Alfonso X, Autónoma y Complutense.},
	{Ayuda a preparación de exámenes como selectividad o MIR.}
}

\cvevent{Nov 2023 - Present}{Profesora}{Academa Inverxo (\href{https://centroinverxo.com/}{centroinverxo.com})}{
	{Profesora de Refuerzo, con especializacion en estudiantes con necesitados educativas.}
}

\cvevent{Feb 2020 - Jun 2022}{Profesora de Ciencias}{Mas Que Clases (\href{https://masqueclases.es/}{masqueclases.es})}{
	{Enseñanza de refuerzo para ciencias a alumnos Bachillerato, ESO y primaria.}
}

\cvevent{2019 - 2020}{Prácticas curriculares}{Universidad Complutense de Madrid (\href{https://www.ucm.es/}{ucm.es})}{
	{Proyecto TFM: Edad al matrimonio y parentesco biológico en matrimonios	consanguíneos en España Central.}
}

\cvevent{Nov 2018 - Jun 2019}{Profesora de Apoyo Escolar}{Academia Torrejón (\href{https://www.facebook.com/academiatorrejon/?locale=es_ES}{fb/academiatorrejon})}{
	{Apoyo escolar de Matemáticas, Física-Química, Geología-Biología y Lengua para primaria, ESO y Bachillerato}
}

\cvevent{Dic 2016 - Jul 2018}{Profesora de Español}{Excellent Global (\textbf{Brasil}) (\href{https://excellentglobal.com.br/}{excellentglobal})}{
	{Profesora de Español en escuela de Idiomas.}
}

\cvevent{Sep 2015 - Nov 2016}{Investigadora}{Instituto de Geología y Paleontología (\textbf{Cuba}) (\href{https://www.igp.minem.cu/}{igp})}{
	{Investigaciones paleontológicas de fósiles de corales y vertebrados superiores.},
	{Parte del equipo que realizó el Catálogo de Foraminíferos y Equinodermos Fósiles de Cuba.},
	{Proyecto con estudiantes de primaria, ESO y bachillerato para fomentar la inserción a carreras relacionadas con las investigaciones paleontológicas.}
}

\cvevent{Jun 2014 - Jun 2015}{Prácticas curriculares}{Instituto Cubano de Antropología (\textbf{Cuba}) (\href{https://www.cips.cu/instituto-cubano-de-antropologia/}{ican})}{
	{Realización de la Tesis de Diploma: Morfometría Geométrica Cráneo-Facial y	Estimación de la Estatura en Poblaciones Aborígenes Pescadores-Recolectores del	Occidente de Cuba.}
}

\cvevent{Jun - Ago 2013}{Prácticas curriculares}{Instituto de Neurología y Neurocirugía (\textbf{Cuba}) (\href{https://instituciones.sld.cu/ineuro/}{ineuro})}{
	{Diagnóstico molecular y seguimiento nutricional en niños con Distrofia muscular de	Duchenne.}
}

%---------------------------------------------------------------------------------------
%	EDUCATION SECTION
%--------------------------------------------------------------------------------------
\newpage
\cvsection{Educación}

\cvevent{Sep 2018 - Jul 2020}{Gruadada de MSc. Máster en Antropología Física y Evolución Humanas}{\href{https://www.uam.es/}{UAM.es}}{
	{El Máster en Antropología Física y Evolución Humana de la UAM ofrece formación sobre la variabilidad, evolución y diversidad genética y ecológica del ser humano, con asignaturas en antropología genética, ecológica, esquelética y forense, combinando teoría evolutiva y práctica de laboratorio y campo.}
}

\cvevent{2010 - 015}{Gruadada de Lic. Biología}{\href{https://www.uh.cu/carrera/biologia/}{uh.cu/Fbio}}{
	{En la Licenciatura en Bilogia de la Universidad de la Habana se ofrece formación sobre campos generales de ciencias biológicas, incluyendo: genética, microbiología, botánica, zoología, ecología, bioquímica, fisiología, biología celular y molecular.}
}

%---------------------------------------------------------------------------------------
%	MATERIALS SECTION
%--------------------------------------------------------------------------------------

\newpage
\cvsection{Materiales Especializados en Educación}

\cvevent{}{Primaria}{}{
	{Lengua},
	{Mates},
	{Inglés},
	{Ciencia}
}

\cvevent{}{ESO}{}{
	{Matemáticas},
	{Física y química},
	{Biología y Geología},
	{Lengua},
	{Apoyo general}
}

\cvevent{}{Universidad}{}{
	{Anatomía},
	{Genética molecular},
	{Genética de poblaciones},
	{Estadística (bioestadística, estadística multivariada)},
	{Osteología},
	{Paleontología},
	{Evolución humana},
	{Antropología física}
}

%============================================================================%
%
%
%
%	DOCUMENT END
%
%
%
%============================================================================%
\end{document}
